% !TeX root = ../教材题库_全部条目_转写工作稿.tex
% Source: XJ-B1-C05-PR-P201-Y0602
% 第5章 / 5.5 三角函数模型的简单应用 / 教材 p.201 / PDF 第208页
% OCR未校对：公式与文字均可能有误；未完成数学TeX转写；禁止直接发给学生
% 原题及相邻原解截图（核对用）：../题目原图/XJ-B1-C05-PR-P201-Y0602-01.webp
% 本文件只在转写完成后作为题干使用；OCR 条目尚不是可靠题干。
\begin{bankitem}{p.201 练习 第2题}{第5章 / 5.5 三角函数模型的简单应用 / 教材 p.201 / PDF 第208页}{XJ-B1-C05-PR-P201-Y0602}{OCR未校对：公式与文字均可能有误；未完成数学TeX转写；禁止直接发给学生}
\begin{ocrdraft}
[教材p.201 / PDF第208页]
\par
2．如图，一个水轮的半径为4m，水轮圆心0距离水面2m，已知水轮每分
\par
钟转动5圈，当水轮上的点P从水中浮现时（图中点PO）开始计算时间．
\par
（D将点P距离水面的高度方(m)表示为时间'（s）的函数；
\par
（2）点P第一次到达最高点大约需要多长时间？
\end{ocrdraft}
\end{bankitem}
